% --- Archon physics formalization source begin ---
% archon:physics
% archon:covers PhyXMiniProblems/problem_phyx_mini_0255.lean
% archon:source-report reports/phyx_mini/problem_phyx_mini_0255.source.json

\chapter{Physics problem phyx\_mini\_0255}
\label{ch:PhyXMiniProblems_problem_phyx_mini_0255}

\paragraph{Problem source.}
\#\# Physical scenario

Figure shows block 1 of mass \$0.200\textbackslash{} kg\$ sliding to the right over a frictionless elevated surface at a speed of \$8.00\textbackslash{} m/s\$. The block undergoes an elastic collision with stationary block 2, which is attached to a spring of spring constant \$1208.5\textbackslash{} N/m\$. (Assume that the spring does not affect the collision.) After the collision, block 2 oscillates in SHM with a period of \$0.140\textbackslash{} s\$, and block 1 slides off the opposite end of the elevated surface, landing a distance \$d\$ from the base of that surface after falling height \$h = 4.90\textbackslash{} m\$.

\#\# Question

What is the value of \$d\$?

\#\# Answer choices

- A: A: 3.0 m

- B: B: 3.5 m

- C: C: 4.0 m

- D: D: 4.5 m

\#\# Figure caption (auxiliary; use the image as primary evidence)

The image depicts a physics-related diagram with the following elements:

- **Objects:**
  - A block labeled "1" in blue, with an arrow indicating it is moving horizontally to the right.
  - A block labeled "2" in orange, positioned to the right of block "1".
  - A spring labeled "k" attached to block "2".

- **Platform:**
  - The blocks and spring are on a horizontal platform.
  - The platform has a step with a vertical drop labeled \textbackslash{}( h \textbackslash{}) and a horizontal distance labeled \textbackslash{}( d \textbackslash{}).

- **Additional Elements:**
  - A dashed curve is shown, suggesting a possible trajectory or path.

The diagram seems to illustrate a mechanical or physics scenario involving motion and forces.

\#\# Dataset metadata

Wave Properties; Multi-Formula Reasoning, Physical Model Grounding Reasoning

\paragraph{Recorded answer/context.}
C: C: 4.0 m

\paragraph{Figure/image path.}
/root/proposal\_for\_physic/science-mango/phyx\_mini\_run/phyx\_data/test\_image/255.png

\paragraph{Formalization target.}
create a compiling Lean file with sorry bodies at `PhyXMiniProblems/problem\_phyx\_mini\_0255.lean`. The Lean declarations must preserve the physical quantities, dimensions or dimensional roles, figure labels, governing-law hypotheses, and final relation expressed by this problem.
Use Mathlib/Physlib names found through LeanExplore where available. If a physics API is missing, introduce faithful local abstractions rather than scalar placeholder aliases.

\begin{theorem}[Physics formalization target]
\label{thm:physics:phyx_mini_0255:target}
\lean{PhyXMiniProblems.ProblemPhyXMini0255.problem_phyx_mini_0255}
\uses{def:physics:phyx-mini-0255:phyxminiproblems-problemphyxmini0255-lengthinmeters, def:physics:phyx-mini-0255:phyxminiproblems-problemphyxmini0255-elasticcollisionspringfallsetup, def:physics:phyx-mini-0255:phyxminiproblems-problemphyxmini0255-matchesprimaryfigure, def:physics:phyx-mini-0255:phyxminiproblems-problemphyxmini0255-matchesproblemdescription, def:physics:phyx-mini-0255:phyxminiproblems-problemphyxmini0255-usesstandardnearearthgravity, def:physics:phyx-mini-0255:phyxminiproblems-problemphyxmini0255-hasphysicalparameters, def:physics:phyx-mini-0255:phyxminiproblems-problemphyxmini0255-satisfiesidealmechanicallaws, def:physics:phyx-mini-0255:phyxminiproblems-problemphyxmini0255-exactlandingdistanceinmeters, def:physics:phyx-mini-0255:phyxminiproblems-problemphyxmini0255-matchesdisplayeddistance}
The assigned autoformalize agent should translate this physics problem into Lean declarations in the covered file, with theorem and lemma proof bodies written as `by sorry`.
\end{theorem}
\begin{proof}
This is an autoformalization task, not a proof task. Produce statements that can later be proved by the physics prover without weakening the physical model.
\end{proof}

% --- Archon declaration topology begin ---
\section{Declaration topology}
\begin{definition}[velocity Dimension]
  \label{def:physics:phyx-mini-0255:phyxminiproblems-problemphyxmini0255-velocitydimension}
  \lean{PhyXMiniProblems.ProblemPhyXMini0255.velocityDimension}
  The physical dimension of velocity, L T⁻¹
\end{definition}
\begin{proof}
  This is the declared model definition of velocity Dimension.
\end{proof}
\begin{definition}[acceleration Dimension]
  \label{def:physics:phyx-mini-0255:phyxminiproblems-problemphyxmini0255-accelerationdimension}
  \lean{PhyXMiniProblems.ProblemPhyXMini0255.accelerationDimension}
  The physical dimension of acceleration, L T⁻²
\end{definition}
\begin{proof}
  This is the declared model definition of acceleration Dimension.
\end{proof}
\begin{definition}[spring Stiffness Dimension]
  \label{def:physics:phyx-mini-0255:phyxminiproblems-problemphyxmini0255-springstiffnessdimension}
  \lean{PhyXMiniProblems.ProblemPhyXMini0255.springStiffnessDimension}
  The physical dimension of linear spring stiffness, M T⁻²
\end{definition}
\begin{proof}
  This is the declared model definition of spring Stiffness Dimension.
\end{proof}
\begin{definition}[Mass Quantity]
  \label{def:physics:phyx-mini-0255:phyxminiproblems-problemphyxmini0255-massquantity}
  \lean{PhyXMiniProblems.ProblemPhyXMini0255.MassQuantity}
  A nonnegative physical mass
\end{definition}
\begin{proof}
  This is the declared model definition of Mass Quantity.
\end{proof}
\begin{definition}[Length Quantity]
  \label{def:physics:phyx-mini-0255:phyxminiproblems-problemphyxmini0255-lengthquantity}
  \lean{PhyXMiniProblems.ProblemPhyXMini0255.LengthQuantity}
  A nonnegative physical length
\end{definition}
\begin{proof}
  This is the declared model definition of Length Quantity.
\end{proof}
\begin{definition}[Time Quantity]
  \label{def:physics:phyx-mini-0255:phyxminiproblems-problemphyxmini0255-timequantity}
  \lean{PhyXMiniProblems.ProblemPhyXMini0255.TimeQuantity}
  A nonnegative physical duration
\end{definition}
\begin{proof}
  This is the declared model definition of Time Quantity.
\end{proof}
\begin{definition}[Velocity Quantity]
  \label{def:physics:phyx-mini-0255:phyxminiproblems-problemphyxmini0255-velocityquantity}
  \lean{PhyXMiniProblems.ProblemPhyXMini0255.VelocityQuantity}
  \uses{def:physics:phyx-mini-0255:phyxminiproblems-problemphyxmini0255-velocitydimension}
  A signed one-dimensional physical velocity
\end{definition}
\begin{proof}
  This is the declared model definition of Velocity Quantity.
\end{proof}
\begin{definition}[Acceleration Quantity]
  \label{def:physics:phyx-mini-0255:phyxminiproblems-problemphyxmini0255-accelerationquantity}
  \lean{PhyXMiniProblems.ProblemPhyXMini0255.AccelerationQuantity}
  \uses{def:physics:phyx-mini-0255:phyxminiproblems-problemphyxmini0255-accelerationdimension}
  A nonnegative magnitude of gravitational acceleration
\end{definition}
\begin{proof}
  This is the declared model definition of Acceleration Quantity.
\end{proof}
\begin{definition}[Spring Stiffness Quantity]
  \label{def:physics:phyx-mini-0255:phyxminiproblems-problemphyxmini0255-springstiffnessquantity}
  \lean{PhyXMiniProblems.ProblemPhyXMini0255.SpringStiffnessQuantity}
  \uses{def:physics:phyx-mini-0255:phyxminiproblems-problemphyxmini0255-springstiffnessdimension}
  A nonnegative physical spring stiffness
\end{definition}
\begin{proof}
  This is the declared model definition of Spring Stiffness Quantity.
\end{proof}
\begin{definition}[nonnegative SIReadout]
  \label{def:physics:phyx-mini-0255:phyxminiproblems-problemphyxmini0255-nonnegativesireadout}
  \lean{PhyXMiniProblems.ProblemPhyXMini0255.nonnegativeSIReadout}
  Read a nonnegative dimensionful quantity in coherent SI units
\end{definition}
\begin{proof}
  This is the declared model definition of nonnegative SIReadout.
\end{proof}
\begin{definition}[signed SIReadout]
  \label{def:physics:phyx-mini-0255:phyxminiproblems-problemphyxmini0255-signedsireadout}
  \lean{PhyXMiniProblems.ProblemPhyXMini0255.signedSIReadout}
  Read a signed dimensionful quantity in coherent SI units
\end{definition}
\begin{proof}
  This is the declared model definition of signed SIReadout.
\end{proof}
\begin{definition}[mass In Kilograms]
  \label{def:physics:phyx-mini-0255:phyxminiproblems-problemphyxmini0255-massinkilograms}
  \lean{PhyXMiniProblems.ProblemPhyXMini0255.massInKilograms}
  \uses{def:physics:phyx-mini-0255:phyxminiproblems-problemphyxmini0255-massquantity, def:physics:phyx-mini-0255:phyxminiproblems-problemphyxmini0255-nonnegativesireadout}
  Kilogram readout of a physical mass
\end{definition}
\begin{proof}
  This is the declared model definition of mass In Kilograms.
\end{proof}
\begin{definition}[length In Meters]
  \label{def:physics:phyx-mini-0255:phyxminiproblems-problemphyxmini0255-lengthinmeters}
  \lean{PhyXMiniProblems.ProblemPhyXMini0255.lengthInMeters}
  \uses{def:physics:phyx-mini-0255:phyxminiproblems-problemphyxmini0255-lengthquantity, def:physics:phyx-mini-0255:phyxminiproblems-problemphyxmini0255-nonnegativesireadout}
  Metre readout of a physical length
\end{definition}
\begin{proof}
  This is the declared model definition of length In Meters.
\end{proof}
\begin{definition}[time In Seconds]
  \label{def:physics:phyx-mini-0255:phyxminiproblems-problemphyxmini0255-timeinseconds}
  \lean{PhyXMiniProblems.ProblemPhyXMini0255.timeInSeconds}
  \uses{def:physics:phyx-mini-0255:phyxminiproblems-problemphyxmini0255-timequantity, def:physics:phyx-mini-0255:phyxminiproblems-problemphyxmini0255-nonnegativesireadout}
  Second readout of a physical duration
\end{definition}
\begin{proof}
  This is the declared model definition of time In Seconds.
\end{proof}
\begin{definition}[velocity In Meters Per Second]
  \label{def:physics:phyx-mini-0255:phyxminiproblems-problemphyxmini0255-velocityinmeterspersecond}
  \lean{PhyXMiniProblems.ProblemPhyXMini0255.velocityInMetersPerSecond}
  \uses{def:physics:phyx-mini-0255:phyxminiproblems-problemphyxmini0255-velocityquantity, def:physics:phyx-mini-0255:phyxminiproblems-problemphyxmini0255-signedsireadout}
  Metre-per-second readout of a signed horizontal velocity
\end{definition}
\begin{proof}
  This is the declared model definition of velocity In Meters Per Second.
\end{proof}
\begin{definition}[acceleration In Meters Per Second Squared]
  \label{def:physics:phyx-mini-0255:phyxminiproblems-problemphyxmini0255-accelerationinmeterspersecondsquared}
  \lean{PhyXMiniProblems.ProblemPhyXMini0255.accelerationInMetersPerSecondSquared}
  \uses{def:physics:phyx-mini-0255:phyxminiproblems-problemphyxmini0255-accelerationquantity, def:physics:phyx-mini-0255:phyxminiproblems-problemphyxmini0255-nonnegativesireadout}
  Metre-per-second-squared readout of an acceleration magnitude
\end{definition}
\begin{proof}
  This is the declared model definition of acceleration In Meters Per Second Squared.
\end{proof}
\begin{definition}[stiffness In Newtons Per Meter]
  \label{def:physics:phyx-mini-0255:phyxminiproblems-problemphyxmini0255-stiffnessinnewtonspermeter}
  \lean{PhyXMiniProblems.ProblemPhyXMini0255.stiffnessInNewtonsPerMeter}
  \uses{def:physics:phyx-mini-0255:phyxminiproblems-problemphyxmini0255-springstiffnessquantity, def:physics:phyx-mini-0255:phyxminiproblems-problemphyxmini0255-nonnegativesireadout}
  Newton-per-metre readout of a spring stiffness
\end{definition}
\begin{proof}
  This is the declared model definition of stiffness In Newtons Per Meter.
\end{proof}
\begin{definition}[Block Label]
  \label{def:physics:phyx-mini-0255:phyxminiproblems-problemphyxmini0255-blocklabel}
  \lean{PhyXMiniProblems.ProblemPhyXMini0255.BlockLabel}
  The two blocks carrying the printed labels 1 and 2
\end{definition}
\begin{proof}
  This is the declared model definition of Block Label.
\end{proof}
\begin{definition}[Horizontal Direction]
  \label{def:physics:phyx-mini-0255:phyxminiproblems-problemphyxmini0255-horizontaldirection}
  \lean{PhyXMiniProblems.ProblemPhyXMini0255.HorizontalDirection}
  Horizontal directions used by the motion arrow and signed velocities
\end{definition}
\begin{proof}
  This is the declared model definition of Horizontal Direction.
\end{proof}
\begin{definition}[Surface Condition]
  \label{def:physics:phyx-mini-0255:phyxminiproblems-problemphyxmini0255-surfacecondition}
  \lean{PhyXMiniProblems.ProblemPhyXMini0255.SurfaceCondition}
  The idealized contact condition stated for the elevated platform
\end{definition}
\begin{proof}
  This is the declared model definition of Surface Condition.
\end{proof}
\begin{definition}[Collision Kind]
  \label{def:physics:phyx-mini-0255:phyxminiproblems-problemphyxmini0255-collisionkind}
  \lean{PhyXMiniProblems.ProblemPhyXMini0255.CollisionKind}
  The type of collision specified in the prose
\end{definition}
\begin{proof}
  This is the declared model definition of Collision Kind.
\end{proof}
\begin{definition}[Spring Influence During Collision]
  \label{def:physics:phyx-mini-0255:phyxminiproblems-problemphyxmini0255-springinfluenceduringcollision}
  \lean{PhyXMiniProblems.ProblemPhyXMini0255.SpringInfluenceDuringCollision}
  How strongly the attached spring acts during the short collision
\end{definition}
\begin{proof}
  This is the declared model definition of Spring Influence During Collision.
\end{proof}
\begin{definition}[Post Collision Motion]
  \label{def:physics:phyx-mini-0255:phyxminiproblems-problemphyxmini0255-postcollisionmotion}
  \lean{PhyXMiniProblems.ProblemPhyXMini0255.PostCollisionMotion}
  The motion observed for block 2 after the collision
\end{definition}
\begin{proof}
  This is the declared model definition of Post Collision Motion.
\end{proof}
\begin{definition}[Figure Element]
  \label{def:physics:phyx-mini-0255:phyxminiproblems-problemphyxmini0255-figureelement}
  \lean{PhyXMiniProblems.ProblemPhyXMini0255.FigureElement}
  Physical objects and points whose relative positions are visible
\end{definition}
\begin{proof}
  This is the declared model definition of Figure Element.
\end{proof}
\begin{definition}[Figure Quantity Label]
  \label{def:physics:phyx-mini-0255:phyxminiproblems-problemphyxmini0255-figurequantitylabel}
  \lean{PhyXMiniProblems.ProblemPhyXMini0255.FigureQuantityLabel}
  The three mathematical labels printed in the primary image
\end{definition}
\begin{proof}
  This is the declared model definition of Figure Quantity Label.
\end{proof}
\begin{definition}[Collision Spring Figure]
  \label{def:physics:phyx-mini-0255:phyxminiproblems-problemphyxmini0255-collisionspringfigure}
  \lean{PhyXMiniProblems.ProblemPhyXMini0255.CollisionSpringFigure}
  \uses{def:physics:phyx-mini-0255:phyxminiproblems-problemphyxmini0255-blocklabel, def:physics:phyx-mini-0255:phyxminiproblems-problemphyxmini0255-horizontaldirection, def:physics:phyx-mini-0255:phyxminiproblems-problemphyxmini0255-figureelement, def:physics:phyx-mini-0255:phyxminiproblems-problemphyxmini0255-figurequantitylabel}
  Categorical and schematic-coordinate data read from the primary image
\end{definition}
\begin{proof}
  This is the declared model definition of Collision Spring Figure.
\end{proof}
\begin{definition}[Elastic Collision Spring Fall Setup]
  \label{def:physics:phyx-mini-0255:phyxminiproblems-problemphyxmini0255-elasticcollisionspringfallsetup}
  \lean{PhyXMiniProblems.ProblemPhyXMini0255.ElasticCollisionSpringFallSetup}
  \uses{def:physics:phyx-mini-0255:phyxminiproblems-problemphyxmini0255-massquantity, def:physics:phyx-mini-0255:phyxminiproblems-problemphyxmini0255-lengthquantity, def:physics:phyx-mini-0255:phyxminiproblems-problemphyxmini0255-timequantity, def:physics:phyx-mini-0255:phyxminiproblems-problemphyxmini0255-velocityquantity, def:physics:phyx-mini-0255:phyxminiproblems-problemphyxmini0255-accelerationquantity, def:physics:phyx-mini-0255:phyxminiproblems-problemphyxmini0255-springstiffnessquantity, def:physics:phyx-mini-0255:phyxminiproblems-problemphyxmini0255-blocklabel, def:physics:phyx-mini-0255:phyxminiproblems-problemphyxmini0255-surfacecondition, def:physics:phyx-mini-0255:phyxminiproblems-problemphyxmini0255-collisionkind, def:physics:phyx-mini-0255:phyxminiproblems-problemphyxmini0255-springinfluenceduringcollision, def:physics:phyx-mini-0255:phyxminiproblems-problemphyxmini0255-postcollisionmotion, def:physics:phyx-mini-0255:phyxminiproblems-problemphyxmini0255-collisionspringfigure}
  The physical quantities for the collision, oscillator, and subsequent fall. The post-collision velocities, flight time, block-2 mass, and landing distance are independent physical quantities. Their relations are supplied only by the governing-law premise below; in particular, the requested value of d is not stored in this structure
\end{definition}
\begin{proof}
  This is the declared model definition of Elastic Collision Spring Fall Setup.
\end{proof}
\begin{definition}[Matches Primary Figure]
  \label{def:physics:phyx-mini-0255:phyxminiproblems-problemphyxmini0255-matchesprimaryfigure}
  \lean{PhyXMiniProblems.ProblemPhyXMini0255.MatchesPrimaryFigure}
  \uses{def:physics:phyx-mini-0255:phyxminiproblems-problemphyxmini0255-elasticcollisionspringfallsetup}
  Primary-image evidence: block 1 lies to the left of block 2 and initially points right; block 2 is attached to the spring and right wall; the opposite platform edge and landing point lie to the left. The image also displays the labels k, h, and d and a dashed projectile path
\end{definition}
\begin{proof}
  This is the declared model definition of Matches Primary Figure.
\end{proof}
\begin{definition}[Matches Problem Description]
  \label{def:physics:phyx-mini-0255:phyxminiproblems-problemphyxmini0255-matchesproblemdescription}
  \lean{PhyXMiniProblems.ProblemPhyXMini0255.MatchesProblemDescription}
  \uses{def:physics:phyx-mini-0255:phyxminiproblems-problemphyxmini0255-massinkilograms, def:physics:phyx-mini-0255:phyxminiproblems-problemphyxmini0255-lengthinmeters, def:physics:phyx-mini-0255:phyxminiproblems-problemphyxmini0255-timeinseconds, def:physics:phyx-mini-0255:phyxminiproblems-problemphyxmini0255-velocityinmeterspersecond, def:physics:phyx-mini-0255:phyxminiproblems-problemphyxmini0255-stiffnessinnewtonspermeter, def:physics:phyx-mini-0255:phyxminiproblems-problemphyxmini0255-elasticcollisionspringfallsetup}
  Numerical readouts and qualitative conditions stated in the problem. The negative post-collision velocity records the prose statement that block 1 leaves from the opposite (left) end, but gives no value for that velocity or for the landing distance
\end{definition}
\begin{proof}
  This is the declared model definition of Matches Problem Description.
\end{proof}
\begin{definition}[Uses Standard Near Earth Gravity]
  \label{def:physics:phyx-mini-0255:phyxminiproblems-problemphyxmini0255-usesstandardnearearthgravity}
  \lean{PhyXMiniProblems.ProblemPhyXMini0255.UsesStandardNearEarthGravity}
  \uses{def:physics:phyx-mini-0255:phyxminiproblems-problemphyxmini0255-accelerationinmeterspersecondsquared, def:physics:phyx-mini-0255:phyxminiproblems-problemphyxmini0255-elasticcollisionspringfallsetup}
  Standard near-Earth gravitational-acceleration readout used by the model
\end{definition}
\begin{proof}
  This is the declared model definition of Uses Standard Near Earth Gravity.
\end{proof}
\begin{definition}[Has Physical Parameters]
  \label{def:physics:phyx-mini-0255:phyxminiproblems-problemphyxmini0255-hasphysicalparameters}
  \lean{PhyXMiniProblems.ProblemPhyXMini0255.HasPhysicalParameters}
  \uses{def:physics:phyx-mini-0255:phyxminiproblems-problemphyxmini0255-massinkilograms, def:physics:phyx-mini-0255:phyxminiproblems-problemphyxmini0255-lengthinmeters, def:physics:phyx-mini-0255:phyxminiproblems-problemphyxmini0255-timeinseconds, def:physics:phyx-mini-0255:phyxminiproblems-problemphyxmini0255-accelerationinmeterspersecondsquared, def:physics:phyx-mini-0255:phyxminiproblems-problemphyxmini0255-stiffnessinnewtonspermeter, def:physics:phyx-mini-0255:phyxminiproblems-problemphyxmini0255-elasticcollisionspringfallsetup}
  Positivity and nondegeneracy of the physical parameters
\end{definition}
\begin{proof}
  This is the declared model definition of Has Physical Parameters.
\end{proof}
\begin{definition}[Satisfies Ideal Mechanical Laws]
  \label{def:physics:phyx-mini-0255:phyxminiproblems-problemphyxmini0255-satisfiesidealmechanicallaws}
  \lean{PhyXMiniProblems.ProblemPhyXMini0255.SatisfiesIdealMechanicalLaws}
  \uses{def:physics:phyx-mini-0255:phyxminiproblems-problemphyxmini0255-massinkilograms, def:physics:phyx-mini-0255:phyxminiproblems-problemphyxmini0255-lengthinmeters, def:physics:phyx-mini-0255:phyxminiproblems-problemphyxmini0255-timeinseconds, def:physics:phyx-mini-0255:phyxminiproblems-problemphyxmini0255-velocityinmeterspersecond, def:physics:phyx-mini-0255:phyxminiproblems-problemphyxmini0255-accelerationinmeterspersecondsquared, def:physics:phyx-mini-0255:phyxminiproblems-problemphyxmini0255-stiffnessinnewtonspermeter, def:physics:phyx-mini-0255:phyxminiproblems-problemphyxmini0255-elasticcollisionspringfallsetup}
  The ideal mechanical model consists of: Physlib's harmonic oscillator for block 2, whose angular frequency is sqrt (k / m) and whose period is 2 pi / omega; one-dimensional momentum and kinetic-energy conservation during the isolated elastic collision; vertical constant-gravity free fall and uniform horizontal motion after block 1 leaves the frictionless platform. These are general governing relations. They contain neither the numerical landing distance nor an answer-choice assertion
\end{definition}
\begin{proof}
  This is the declared model definition of Satisfies Ideal Mechanical Laws.
\end{proof}
\begin{definition}[exact Landing Distance In Meters]
  \label{def:physics:phyx-mini-0255:phyxminiproblems-problemphyxmini0255-exactlandingdistanceinmeters}
  \lean{PhyXMiniProblems.ProblemPhyXMini0255.exactLandingDistanceInMeters}
  The exact SI expression obtained from the measured spring period, the one-dimensional elastic-collision formula, and horizontal projectile motion. This is only a scalar expression built from the stated data; the independent physical field setup.landingDistance is not defined to equal it
\end{definition}
\begin{proof}
  This is the declared model definition of exact Landing Distance In Meters.
\end{proof}
\begin{definition}[Answer Choice]
  \label{def:physics:phyx-mini-0255:phyxminiproblems-problemphyxmini0255-answerchoice}
  \lean{PhyXMiniProblems.ProblemPhyXMini0255.AnswerChoice}
  Labels of the four numerical choices printed with the problem
\end{definition}
\begin{proof}
  This is the declared model definition of Answer Choice.
\end{proof}
\begin{definition}[distance In Meters]
  \label{def:physics:phyx-mini-0255:phyxminiproblems-problemphyxmini0255-answerchoice-distanceinmeters}
  \lean{PhyXMiniProblems.ProblemPhyXMini0255.AnswerChoice.distanceInMeters}
  \uses{def:physics:phyx-mini-0255:phyxminiproblems-problemphyxmini0255-answerchoice}
  Distance in metres printed beside each displayed answer choice
\end{definition}
\begin{proof}
  This is the declared model definition of distance In Meters.
\end{proof}
\begin{definition}[recorded Answer Choice]
  \label{def:physics:phyx-mini-0255:phyxminiproblems-problemphyxmini0255-recordedanswerchoice}
  \lean{PhyXMiniProblems.ProblemPhyXMini0255.recordedAnswerChoice}
  \uses{def:physics:phyx-mini-0255:phyxminiproblems-problemphyxmini0255-answerchoice}
  The answer label recorded by the supplied dataset metadata
\end{definition}
\begin{proof}
  This is the declared model definition of recorded Answer Choice.
\end{proof}
\begin{definition}[Matches Displayed Distance]
  \label{def:physics:phyx-mini-0255:phyxminiproblems-problemphyxmini0255-matchesdisplayeddistance}
  \lean{PhyXMiniProblems.ProblemPhyXMini0255.MatchesDisplayedDistance}
  \uses{def:physics:phyx-mini-0255:phyxminiproblems-problemphyxmini0255-lengthquantity, def:physics:phyx-mini-0255:phyxminiproblems-problemphyxmini0255-lengthinmeters, def:physics:phyx-mini-0255:phyxminiproblems-problemphyxmini0255-answerchoice}
  Agreement of the physical distance with a displayed nearest-tenth metre
\end{definition}
\begin{proof}
  This is the declared model definition of Matches Displayed Distance.
\end{proof}
\begin{lemma}[landing Distance eq exact Expression]
  \label{lem:physics:phyx-mini-0255:phyxminiproblems-problemphyxmini0255-landingdistance-eq-exactexpression}
  \lean{PhyXMiniProblems.ProblemPhyXMini0255.landingDistance_eq_exactExpression}
  \uses{def:physics:phyx-mini-0255:phyxminiproblems-problemphyxmini0255-lengthinmeters, def:physics:phyx-mini-0255:phyxminiproblems-problemphyxmini0255-elasticcollisionspringfallsetup, def:physics:phyx-mini-0255:phyxminiproblems-problemphyxmini0255-matchesprimaryfigure, def:physics:phyx-mini-0255:phyxminiproblems-problemphyxmini0255-matchesproblemdescription, def:physics:phyx-mini-0255:phyxminiproblems-problemphyxmini0255-usesstandardnearearthgravity, def:physics:phyx-mini-0255:phyxminiproblems-problemphyxmini0255-hasphysicalparameters, def:physics:phyx-mini-0255:phyxminiproblems-problemphyxmini0255-satisfiesidealmechanicallaws, def:physics:phyx-mini-0255:phyxminiproblems-problemphyxmini0255-exactlandingdistanceinmeters}
  The governing laws determine the exact (unrounded) landing-distance readout
\end{lemma}
\begin{proof}
  The conclusion follows from the governing relations and figure readouts named in the statement of landing Distance eq exact Expression.
\end{proof}
\begin{lemma}[exact Landing Distance matches choice C]
  \label{lem:physics:phyx-mini-0255:phyxminiproblems-problemphyxmini0255-exactlandingdistance-matches-choice-c}
  \lean{PhyXMiniProblems.ProblemPhyXMini0255.exactLandingDistance_matches_choice_C}
  \uses{def:physics:phyx-mini-0255:phyxminiproblems-problemphyxmini0255-exactlandingdistanceinmeters, def:physics:phyx-mini-0255:phyxminiproblems-problemphyxmini0255-answerchoice}
  The exact expression is approximately 3.99994 m, so it agrees with the displayed nearest-tenth value 4.0 m
\end{lemma}
\begin{proof}
  The conclusion follows from the governing relations and figure readouts named in the statement of exact Landing Distance matches choice C.
\end{proof}
% --- Archon declaration topology end ---
% --- Archon physics formalization source end ---
